\documentclass[10pt,a4paper]{amsart}
\usepackage{amsmath,amssymb,amsthm,graphicx}
\usepackage{verbatim}
\newenvironment{reference}{\comment}{\endcomment}
%\newenvironment{reference}{}{}
\newenvironment{slogan}{\comment}{\endcomment}
\newenvironment{history}{\comment}{\endcomment}

% For commutative diagrams we use Xy-pic
\usepackage[all]{xy}

% We use 2cell for 2-commutative diagrams.
\xyoption{2cell}
% Theorem environments.
%
\theoremstyle{plain}
\newtheorem{theorem}[subsection]{Theorem}
\newtheorem{proposition}[subsection]{Proposition}
\newtheorem{lemma}[subsection]{Lemma}

\theoremstyle{definition}
\newtheorem{definition}[subsection]{Definition}
\newtheorem{example}[subsection]{Example}
\newtheorem{exercise}[subsection]{Exercise}
\newtheorem{situation}[subsection]{Situation}

\theoremstyle{remark}
\newtheorem{remark}[subsection]{Remark}
\newtheorem{remarks}[subsection]{Remarks}

\numberwithin{equation}{subsection}

% Macros
%
\def\lim{\mathop{\mathrm{lim}}\nolimits}
\def\colim{\mathop{\mathrm{colim}}\nolimits}
\def\Spec{\mathop{\mathrm{Spec}}}
\def\Hom{\mathop{\mathrm{Hom}}\nolimits}
\def\Ext{\mathop{\mathrm{Ext}}\nolimits}
\def\SheafHom{\mathop{\mathcal{H}\!\mathit{om}}\nolimits}
\def\SheafExt{\mathop{\mathcal{E}\!\mathit{xt}}\nolimits}
\def\Sch{\mathit{Sch}}
\def\Mor{\mathop{\mathrm{Mor}}\nolimits}
\def\Ob{\mathop{\mathrm{Ob}}\nolimits}
\def\Sh{\mathop{\mathit{Sh}}\nolimits}
\def\NL{\mathop{N\!L}\nolimits}
\def\CH{\mathop{\mathrm{CH}}\nolimits}
\def\proetale{{pro\text{-}\acute{e}tale}}
\def\etale{{\acute{e}tale}}
\def\QCoh{\mathit{QCoh}}
\def\Ker{\mathop{\mathrm{Ker}}}
\def\Im{\mathop{\mathrm{Im}}}
\def\Coker{\mathop{\mathrm{Coker}}}
\def\Coim{\mathop{\mathrm{Coim}}}

% Boxtimes
%
\DeclareMathSymbol{\boxtimes}{\mathbin}{AMSa}{"02}

%
% Macros for moduli stacks/spaces
%
\def\QCohstack{\mathcal{QC}\!\mathit{oh}}
\def\Cohstack{\mathcal{C}\!\mathit{oh}}
\def\Spacesstack{\mathcal{S}\!\mathit{paces}}
\def\Quotfunctor{\mathrm{Quot}}
\def\Hilbfunctor{\mathrm{Hilb}}
\def\Curvesstack{\mathcal{C}\!\mathit{urves}}
\def\Polarizedstack{\mathcal{P}\!\mathit{olarized}}
\def\Complexesstack{\mathcal{C}\!\mathit{omplexes}}
% \Pic is the operator that assigns to X its picard group, usage \Pic(X)
% \Picardstack_{X/B} denotes the Picard stack of X over B
% \Picardfunctor_{X/B} denotes the Picard functor of X over B
\def\Pic{\mathop{\mathrm{Pic}}\nolimits}
\def\Picardstack{\mathcal{P}\!\mathit{ic}}
\def\Picardfunctor{\mathrm{Pic}}
\def\Deformationcategory{\mathcal{D}\!\mathit{ef}}
\usepackage{cleveref}
\crefname{originallemma}{original lemma}{original lemmas}
\Crefname{originallemma}{Original lemma}{Original lemmas}
\makeatletter
\newlabel{lemma-precompose-property-local-source}{{06EQ}{source}{}{sourcepointer-1}{}}
\newlabel{lemma-precompose-property-local-source@cref}{{[originallemma][0][]06EQ}{[1][0][]source}{}{}{}}
\makeatother
\crefname{originallemma}{original lemma}{original lemmas}
\Crefname{originallemma}{Original lemma}{Original lemmas}
\makeatletter
\newlabel{lemma-pullback-property-local-target}{{06EM}{source}{}{sourcepointer-2}{}}
\newlabel{lemma-pullback-property-local-target@cref}{{[originallemma][0][]06EM}{[1][0][]source}{}{}{}}
\makeatother
\crefname{originaldefinition}{original definition}{original definitions}
\Crefname{originaldefinition}{Original definition}{Original definitions}
\makeatletter
\newlabel{definition-local-source-target}{{06F8}{source}{}{sourcepointer-3}{}}
\newlabel{definition-local-source-target@cref}{{[originaldefinition][0][]06F8}{[1][0][]source}{}{}{}}
\makeatother
\crefname{originalreference}{original reference}{original references}
\Crefname{originalreference}{Original reference}{Original references}
\makeatletter
\newlabel{descent-lemma-etale-tau-local-source-target}{{0CF1}{source}{}{sourcepointer-4}{}}
\newlabel{descent-lemma-etale-tau-local-source-target@cref}{{[originalreference][0][]0CF1}{[1][0][]source}{}{}{}}
\makeatother
\title{\textbackslash{}'Etale-smooth local on the source-and-target from three local properties}
\author{The Stacks Project Authors}
\date{stacks-project:spaces-descent.tex:local-snapshot-20261012; GNU FDL 1.2 or later}
\begin{document}
\maketitle
\noindent\textit{Editorial note.} All delivered mathematical text is the author's own, copied verbatim from the pinned Stacks Project chapter file spaces-descent.tex; no statement or proof step was rewritten, reordered or removed.\par\medskip
\noindent\textit{Editorial note.} The pinned source reads 'with \$X\_i\$ a scheme' at proof line 3367 where \$X'\$ is meant; delivered verbatim per policy, flagged here for the reviewer.\par\medskip
\noindent\textit{Allowed standard prerequisites (editorial).}\begin{itemize}
\item Definition of \textbackslash{}'etale-smooth local on the source-and-target. Stacks Project, Descent on Spaces, Definition definition-local-source-target. Supplies the target notion, in particular part (3)(b) whose data the proof assembles.
\item Stability under \textbackslash{}'etale precomposition. Stacks Project, Descent on Spaces, Lemma lemma-precompose-property-local-source. Derives stability under precomposing with \textbackslash{}'etale morphisms from assumption (1).
\item Stability under smooth base change. Stacks Project, Descent on Spaces, Lemma lemma-pullback-property-local-target. Derives stability under smooth base change from assumption (2).
\item Scheme-side \textbackslash{}'etale-smooth criterion. Stacks Project, Descent, Lemma descent-lemma-etale-tau-local-source-target. Concludes for morphisms of schemes after the reduction.
\end{itemize}\medskip
\ifdefined\hypertarget\hypertarget{sourcepointer-1}{}\fi
\noindent\textit{Source reference.} Original lemma 06EQ: sources/stacks-chow/tags.txt, source lines 8224--8224. Cross-chapter standard prerequisite cited by the author inside the counted span; it is used as a pointer and its own proof is not part of this item.\par\medskip
\ifdefined\hypertarget\hypertarget{sourcepointer-2}{}\fi
\noindent\textit{Source reference.} Original lemma 06EM: sources/stacks-chow/tags.txt, source lines 8221--8221. Cross-chapter standard prerequisite cited by the author inside the counted span; it is used as a pointer and its own proof is not part of this item.\par\medskip
\ifdefined\hypertarget\hypertarget{sourcepointer-3}{}\fi
\noindent\textit{Source reference.} Original definition 06F8: sources/stacks-chow/tags.txt, source lines 8242--8242. Cross-chapter standard prerequisite cited by the author inside the counted span; it is used as a pointer and its own proof is not part of this item.\par\medskip
\ifdefined\hypertarget\hypertarget{sourcepointer-4}{}\fi
\noindent\textit{Source reference.} Original reference 0CF1: sources/stacks-chow/tags.txt, source lines 15596--15596. Cross-chapter standard prerequisite cited by the author inside the counted span; it is used as a pointer and its own proof is not part of this item.\par\medskip
\begin{lemma}
\label{lemma-etale-smooth-local-source-target}
Let $S$ be a scheme.
Let $\mathcal{P}$ be a property of morphisms of algebraic spaces over $S$.
Assume
\begin{enumerate}
\item $\mathcal{P}$ is \'etale local on the source,
\item $\mathcal{P}$ is smooth local on the target, and
\item $\mathcal{P}$ is stable under postcomposing with open immersions:
if $f : X \to Y$ has $\mathcal{P}$ and $Y \subset Z$ is an open embedding
then $X \to Z$ has $\mathcal{P}$.
\end{enumerate}
Then $\mathcal{P}$ is \'etale-smooth local on the source-and-target.
\end{lemma}

\begin{proof}
Let $\mathcal{P}$ be a property of morphisms of algebraic spaces which
satisfies conditions (1), (2) and (3) of the lemma. By
Lemma \ref{lemma-precompose-property-local-source}
we see that $\mathcal{P}$ is stable under precomposing with
\'etale morphisms. By
Lemma \ref{lemma-pullback-property-local-target}
we see that $\mathcal{P}$ is stable under smooth base change.
Hence it suffices to prove part (3) of
Definition \ref{definition-local-source-target}
holds.

\medskip\noindent
More precisely, suppose that $f : X \to Y$ is a morphism
of algebraic spaces over $S$ which satisfies
Definition \ref{definition-local-source-target} part (3)(b).
In other words, for every $x \in X$ there exists
a smooth morphism $b_x : V_x \to Y$,
an \'etale morphism $U_x \to V_x \times_Y X$, and
a point $u_x \in |U_x|$ mapping to $x$
such that $h_x : U_x \to V_x$ has $\mathcal{P}$.
The proof of the lemma is complete once we show that $f$ has $\mathcal{P}$.

\medskip\noindent
Let $a_x : U_x \to X$ be the composition $U_x \to V_x \times_Y X \to X$.
Set $U = \coprod U_x$, $a = \coprod a_x$, $V = \coprod V_x$,
$b = \coprod b_x$, and $h = \coprod h_x$. We obtain a
commutative diagram
$$
\xymatrix{
U \ar[d]_a \ar[r]_h & V \ar[d]^b \\
X \ar[r]^f & Y
}
$$
with $b$ smooth, $U \to V \times_Y X$ \'etale, $a$ surjective.
Note that $h$ has $\mathcal{P}$ as each $h_x$ does and $\mathcal{P}$
is smooth local on the target. In the next paragraph we prove
that we may assume $U, V, X, Y$ are schemes; we encourage the reader
to skip it.

\medskip\noindent
Let $X, Y, U, V, a, b, f, h$ be as in the previous paragraph. We have
to show $f$ has $\mathcal{P}$. Let $X' \to X$ be a surjective \'etale
morphism with $X_i$ a scheme. Set $U' = X' \times_X U$. Then
$U' \to X'$ is surjective and $U' \to X' \times_Y V$ is \'etale.
Since $\mathcal{P}$ is \'etale local on the source, we see that
$U' \to V$ has $\mathcal{P}$ and that it suffices to show that
$X' \to Y$ has $\mathcal{P}$. In other words, we may assume
that $X$ is a scheme. Next, choose a surjective \'etale morphism
$Y' \to Y$ with $Y'$ a scheme. Set $V' = V \times_Y Y'$,
$X' = X \times_Y Y'$, and $U' = U \times_Y Y'$. Then
$U' \to X'$ is surjective and $U' \to X' \times_{Y'} V'$ is \'etale.
Since $\mathcal{P}$ is smooth local on the target, we see that $U' \to V'$ has
$\mathcal{P}$ and that it suffices to prove $X' \to Y'$ has $\mathcal{P}$.
Thus we may assume both $X$ and $Y$ are schemes.
Choose a surjective \'etale morphism $V' \to V$
with $V'$ a scheme. Set $U' = U \times_V V'$.
Then $U' \to X$ is surjective and $U' \to X \times_Y V'$ is \'etale.
Since $\mathcal{P}$ is smooth local on the source, we see that
$U' \to V'$ has $\mathcal{P}$. Thus we may replace $U, V$ by
$U', V'$ and assume $X, Y, V$ are schemes.
Finally, we replace $U$ by a scheme surjective \'etale over $U$
and we see that we may assume $U, V, X, Y$ are all schemes.

\medskip\noindent
If $U, V, X, Y$ are schemes, then $f$ has $\mathcal{P}$
by Descent, Lemma \ref{descent-lemma-etale-tau-local-source-target}.
\end{proof}

\end{document}
